\documentclass[hidelinks]{article}
\usepackage[letterpaper, total={7.5in, 10in}]{geometry} % Paper size
\usepackage[dvipsnames]{xcolor} %colors
\usepackage{graphicx} % Required for inserting images
\usepackage{amsmath} % Math formatting, ex. alignment, and better symbols, idk how this really works 
\usepackage{tikz}  % Tikz for drawing
\usepackage{tkz-euclide} 
\usepackage[unicode]{hyperref} % For hyperlinks and cross referencing
\usepackage[all]{hypcap} 
\usepackage{fancyhdr} % For fancy headers
\usepackage{comment} % For comments in the code

\usetikzlibrary{angles,calc, decorations.pathreplacing}

\definecolor{carminered}{rgb}{1.0, 0.0, 0.22}
\definecolor{capri}{rgb}{0.0, 0.75, 1.0}
\definecolor{brightlavender}{rgb}{0.75, 0.58, 0.89}

\title{\textbf{PHYS-130-1411 Homework 1}} % Title of the document
\author{Ziyang Ling}
\date{\today} % Date of the document
\begin{document}
\hypersetup{bookmarksnumbered=true,}
\pagecolor{white}
\color{black}
\maketitle % Prints the title information we gave before, ex. title, date, author

% \begin{Large}
% \tableofcontents
% \end{Large}%
% \pagebreak

% For table of contents, uncomment the above lines


\section{Question 1}
Question text

{\color{SkyBlue}
Answer work

\color{Thistle}{\textbf{Answer}}
}

\end{document}



\begin{comment}
    
\begin{tikzpicture}

% coords for drawing angles
\coordinate (origin) at (0,0);
\coordinate (negx) at (-5,0);
\coordinate (posx) at (5,0);
\coordinate (negy) at (0, -5);
\coordinate (vecb) at (2.4075, -9.1945);
\coordinate (veca) at (-2.4075, 3.1945);

%axes
\draw[<->,ultra thick] (-5,0) coordinate (A) --(5,0) node[right]{$x$};
\draw[<->,ultra thick] (0,-5)--(0,5) node[above]{$y$};

%vec a
\draw[->, thick, capri] (0,0) coordinate (B) -- (-2.4075, 3.1945) coordinate (C) node[left]{\textbf{A}};

% angle for vec a
\draw pic[thick, capri, "53$^\circ$", draw, <-,angle radius=1cm,angle eccentricity=1.4] {angle = veca--origin--negx};

% opposite angle for vec a
\draw pic[thick, "127$^\circ$", draw, ->,angle radius=1cm, angle eccentricity=1.4] {angle = posx--origin--veca};

% magnitude label for vec a
\draw [thick, capri, decorate,decoration={brace,amplitude=6pt,mirror,raise=1ex}] (0,0) -- (-2.4075, 3.1945) node[midway, yshift=2em, xshift=2em]{8.0 m};

% vec b
\draw[->, thick, carminered] (0,0) coordinate (B) -- (2.4075, -9.1945) coordinate (C) node[left]{\textbf{B}};

% angle for vec b
\draw pic[thick, carminered, "14.7$^\circ$", draw, ->,angle radius=5cm, angle eccentricity=1.1] {angle = negy--origin--vecb};

% label for vec b
\draw [thick, carminered, decorate,decoration={brace,amplitude=6pt,raise=1ex}] (0,0) -- (2.4075, -9.1945) node[midway, xshift=2.5em, yshift=0.5em]{19.0 m};

%vec c
\draw[->, thick, brightlavender] (0,0) coordinate (B) -- (0, -6) coordinate (C) node[below]{\textbf{A + B}};

% label for vec c
\draw [thick, brightlavender, decorate,decoration={brace,amplitude=6pt,raise=1ex, mirror}] (0,0) -- (0, -6) node[midway, xshift=-3em]{12.0 m};
\end{tikzpicture}

\end{comment}
%Uncomment the above code to draw a vector diagram